\documentclass[12pt, a4paper]{report}

% ============================================================
% PACKAGES
% ============================================================
\usepackage[utf8]{inputenc}
\usepackage[T1]{fontenc}
\usepackage{amsmath, amssymb, amsthm}
\usepackage{mathtools}
\usepackage{geometry}
\usepackage{graphicx}
\usepackage{hyperref}
\usepackage{xcolor}
\usepackage{tcolorbox}
\usepackage{listings}
\usepackage{booktabs}
\usepackage{tikz}
\usepackage{pgfplots}
\usepackage{fancyhdr}
\usepackage{titlesec}
\usepackage{enumitem}
\usepackage{float}

\usetikzlibrary{arrows.meta, positioning, calc, decorations.markings}
\pgfplotsset{compat=1.18}

\geometry{margin=1in}
\hypersetup{
    colorlinks=true,
    linkcolor=blue!70!black,
    urlcolor=blue!60!black,
    citecolor=green!50!black
}

% ============================================================
% THEOREM ENVIRONMENTS
% ============================================================
\theoremstyle{definition}
\newtheorem{definition}{Definition}[chapter]
\newtheorem{example}{Example}[chapter]
\newtheorem{theorem}{Theorem}[chapter]
\newtheorem{lemma}[theorem]{Lemma}

% ============================================================
% CUSTOM BOXES
% ============================================================
\newtcolorbox{keyidea}[1][]{
    colback=blue!5!white,
    colframe=blue!75!black,
    fonttitle=\bfseries,
    title={Key Idea},
    #1
}

\newtcolorbox{numexample}[1][]{
    colback=green!5!white,
    colframe=green!50!black,
    fonttitle=\bfseries,
    title={Numerical Example},
    #1
}

\newtcolorbox{warning}[1][]{
    colback=red!5!white,
    colframe=red!75!black,
    fonttitle=\bfseries,
    title={Common Mistake},
    #1
}

% ============================================================
% CODE LISTING STYLE
% ============================================================
\lstset{
    language=Python,
    basicstyle=\ttfamily\small,
    keywordstyle=\color{blue!80!black}\bfseries,
    commentstyle=\color{green!50!black}\itshape,
    stringstyle=\color{red!60!black},
    showstringspaces=false,
    frame=single,
    backgroundcolor=\color{gray!10},
    breaklines=true,
    numbers=left,
    numberstyle=\tiny\color{gray}
}

% ============================================================
% HEADERS / FOOTERS
% ============================================================
\pagestyle{fancy}
\fancyhf{}
\fancyhead[L]{\leftmark}
\fancyhead[R]{Mini-Keras}
\fancyfoot[C]{\thepage}

% ============================================================
% TITLE
% ============================================================
\title{
    \vspace{-1cm}
    {\Huge\bfseries Building a Neural Network Framework\\[0.3cm] from Scratch}\\[1cm]
    {\Large The Mathematics Behind Mini-Keras}\\[0.5cm]
    {\large A Beginner's Guide with Proofs \& Numerical Examples}
}
\author{
    \textbf{Abel Yohannes}\\
    \texttt{https://github.com/abelyo252/}
}
\date{\today}

% ============================================================
\begin{document}
\maketitle
\tableofcontents
\newpage

% ============================================================
% CHAPTER 1: INTRODUCTION
% ============================================================
\chapter{Introduction}

\section{What Is a Neural Network?}

A neural network is a mathematical function $f: \mathbb{R}^n \to \mathbb{R}^m$ that maps an input vector $\mathbf{x} \in \mathbb{R}^n$ to an output vector $\hat{\mathbf{y}} \in \mathbb{R}^m$. It is composed of a sequence of simpler functions called \textit{layers}, each involving a linear transformation followed by a non-linear \textit{activation function}.

\begin{definition}[Artificial Neuron]
A single neuron computes:
\begin{equation}
    z = \sum_{i=1}^{n} w_i x_i + b = \mathbf{w}^\top \mathbf{x} + b
\end{equation}
where $\mathbf{w} \in \mathbb{R}^n$ is the weight vector, $b \in \mathbb{R}$ is the bias, and $\mathbf{x} \in \mathbb{R}^n$ is the input vector. The output of the neuron is $a = \sigma(z)$, where $\sigma$ is an activation function.
\end{definition}

\begin{numexample}
Consider a single neuron with 2 inputs:
\[
\mathbf{x} = \begin{bmatrix} 0.8 & -0.5 \end{bmatrix}, \quad
\mathbf{w} = \begin{bmatrix} 1.2 & 0.4 \end{bmatrix}, \quad
b = -0.1
\]
\textbf{Step 1 (Linear combination)}:
\[
z = w_1 x_1 + w_2 x_2 + b = (1.2)(0.8) + (0.4)(-0.5) - 0.1 = 0.96 - 0.20 - 0.10 = \boxed{0.66}
\]
\textbf{Step 2 (Activation via Sigmoid)}:
\[
a = \sigma(z) = \frac{1}{1 + e^{-0.66}} = \frac{1}{1 + 0.5169} = \boxed{0.6593}
\]
\end{numexample}

\section{Architecture of a Feedforward Network}

A feedforward neural network with $L$ layers can be written as:
\begin{equation}
    \hat{\mathbf{y}} = f(\mathbf{x}) = \sigma_L\bigl(W_L \cdot \sigma_{L-1}\bigl(\cdots \sigma_1(W_1 \mathbf{x} + \mathbf{b}_1) \cdots\bigr) + \mathbf{b}_L\bigr)
\end{equation}

For each layer $\ell = 1, 2, \ldots, L$:
\begin{align}
    \mathbf{z}^{[\ell]} &= W^{[\ell]} \mathbf{a}^{[\ell-1]} + \mathbf{b}^{[\ell]} \quad \text{(linear transformation)} \\
    \mathbf{a}^{[\ell]} &= \sigma_\ell(\mathbf{z}^{[\ell]}) \quad \text{(activation)}
\end{align}
where $\mathbf{a}^{[0]} = \mathbf{x}$ (the input).

\begin{figure}[H]
\centering
\begin{tikzpicture}[
    node distance=1.5cm,
    neuron/.style={circle, draw, minimum size=0.8cm, inner sep=0pt},
    layer/.style={rectangle, draw, dashed, inner sep=0.3cm}
]
    % Input layer
    \node[neuron, fill=blue!20] (x1) at (0,2) {$x_1$};
    \node[neuron, fill=blue!20] (x2) at (0,0) {$x_2$};
    \node[neuron, fill=blue!20] (x3) at (0,-2) {$x_3$};

    % Hidden layer
    \node[neuron, fill=orange!30] (h1) at (3,1.5) {$a_1^{[1]}$};
    \node[neuron, fill=orange!30] (h2) at (3,0) {$a_2^{[1]}$};
    \node[neuron, fill=orange!30] (h3) at (3,-1.5) {$a_3^{[1]}$};

    % Output layer
    \node[neuron, fill=green!30] (o1) at (6,0.75) {$a_1^{[2]}$};
    \node[neuron, fill=green!30] (o2) at (6,-0.75) {$a_2^{[2]}$};

    % Connections input -> hidden
    \foreach \i in {1,2,3}
        \foreach \j in {1,2,3}
            \draw[->, gray!60] (x\i) -- (h\j);

    % Connections hidden -> output
    \foreach \i in {1,2,3}
        \foreach \j in {1,2}
            \draw[->, gray!60] (h\i) -- (o\j);

    % Labels
    \node[above=0.3cm] at (0,2.5) {\footnotesize Input Layer};
    \node[above=0.3cm] at (3,2.0) {\footnotesize Hidden Layer};
    \node[above=0.3cm] at (6,1.3) {\footnotesize Output Layer};
\end{tikzpicture}
\caption{A feedforward neural network with one hidden layer.}
\end{figure}

\section{The Learning Problem}

\textbf{Goal}: Find parameters $\theta = \{W^{[\ell]}, \mathbf{b}^{[\ell]}\}_{\ell=1}^{L}$ that minimize a loss function:
\begin{equation}
    \theta^* = \arg\min_\theta \frac{1}{N} \sum_{i=1}^{N} \mathcal{L}(\hat{\mathbf{y}}_i, \mathbf{y}_i)
\end{equation}
where $N$ is the number of training samples, $\hat{\mathbf{y}}_i = f(\mathbf{x}_i; \theta)$ is the prediction, and $\mathbf{y}_i$ is the ground truth label.


% ============================================================
% CHAPTER 2: FORWARD PROPAGATION
% ============================================================
\chapter{Forward Propagation}

\section{The Dense (Fully-Connected) Layer}

\begin{definition}[Dense Layer]
Given an input matrix $X \in \mathbb{R}^{N \times d_\text{in}}$ (a batch of $N$ samples with $d_\text{in}$ features), a Dense layer computes:
\begin{equation}
    Z = X W + \mathbf{b}
\end{equation}
where $W \in \mathbb{R}^{d_\text{in} \times d_\text{out}}$ is the weight matrix and $\mathbf{b} \in \mathbb{R}^{d_\text{out}}$ is the bias vector (broadcast across all $N$ samples).
\end{definition}

\begin{numexample}
Let $d_\text{in} = 3$, $d_\text{out} = 2$, and batch size $N = 1$. Consider:
\[
\mathbf{x} = \begin{bmatrix} 1.0 & 2.0 & 3.0 \end{bmatrix}, \quad
W = \begin{bmatrix} 0.1 & 0.4 \\ 0.2 & 0.5 \\ 0.3 & 0.6 \end{bmatrix}, \quad
\mathbf{b} = \begin{bmatrix} 0.1 & 0.2 \end{bmatrix}
\]

\textbf{Step 1}: Compute $\mathbf{z} = \mathbf{x} W + \mathbf{b}$:
\begin{align*}
z_1 &= (1.0)(0.1) + (2.0)(0.2) + (3.0)(0.3) + 0.1 = 0.1 + 0.4 + 0.9 + 0.1 = \boxed{1.5} \\
z_2 &= (1.0)(0.4) + (2.0)(0.5) + (3.0)(0.6) + 0.2 = 0.4 + 1.0 + 1.8 + 0.2 = \boxed{3.4}
\end{align*}

So $\mathbf{z} = \begin{bmatrix} 1.5 & 3.4 \end{bmatrix}$.
\end{numexample}

\begin{numexample}[Batch Processing with $N=2$ Samples]
Now consider a batch of $N=2$ samples with $d_\text{in} = 2$ features and $d_\text{out} = 3$ output units:
\[
X = \begin{bmatrix} 1.0 & 0.5 \\ -0.5 & 2.0 \end{bmatrix}, \quad
W = \begin{bmatrix} 0.2 & -0.1 & 0.4 \\ 0.5 & 0.3 & -0.2 \end{bmatrix}, \quad
\mathbf{b} = \begin{bmatrix} 0.1 & 0.0 & -0.1 \end{bmatrix}
\]

\textbf{Step 1: Compute matrix multiplication $X W$}:
\[
X W = \begin{bmatrix} 
(1.0)(0.2) + (0.5)(0.5) & (1.0)(-0.1) + (0.5)(0.3) & (1.0)(0.4) + (0.5)(-0.2) \\
(-0.5)(0.2) + (2.0)(0.5) & (-0.5)(-0.1) + (2.0)(0.3) & (-0.5)(0.4) + (2.0)(-0.2)
\end{bmatrix}
\]
\[
X W = \begin{bmatrix} 
0.2 + 0.25 & -0.1 + 0.15 & 0.4 - 0.1 \\
-0.1 + 1.0 & 0.05 + 0.6 & -0.2 - 0.4
\end{bmatrix} = \begin{bmatrix} 0.45 & 0.05 & 0.30 \\ 0.90 & 0.65 & -0.60 \end{bmatrix}
\]

\textbf{Step 2: Add bias vector $\mathbf{b}$ (broadcast across rows)}:
\[
Z = X W + \mathbf{b} = \begin{bmatrix} 0.45+0.1 & 0.05+0.0 & 0.30-0.1 \\ 0.90+0.1 & 0.65+0.0 & -0.60-0.1 \end{bmatrix} = \boxed{\begin{bmatrix} 0.55 & 0.05 & 0.20 \\ 1.00 & 0.65 & -0.70 \end{bmatrix}}
\]
\end{numexample}

\section{Weight Initialization}

Proper initialization prevents vanishing or exploding gradients in deep networks.

\begin{definition}[He Initialization]
For layers with ReLU activation:
\begin{equation}
    W_{ij} \sim \mathcal{N}\left(0,\; \sqrt{\frac{2}{d_\text{in}}}\right)
\end{equation}
\end{definition}

\begin{definition}[Xavier (Glorot) Initialization]
For layers with symmetric activations (Sigmoid, Tanh):
\begin{equation}
    W_{ij} \sim \mathcal{U}\left(-\sqrt{\frac{6}{d_\text{in} + d_\text{out}}},\; \sqrt{\frac{6}{d_\text{in} + d_\text{out}}}\right)
\end{equation}
\end{definition}

\begin{numexample}[Weight Initialization Bounds]
Suppose a layer maps $d_\text{in} = 784$ features to $d_\text{out} = 128$ neurons:
\begin{enumerate}
    \item \textbf{He Initialization (for ReLU)}:
    \[
    \sigma = \sqrt{\frac{2}{784}} = \sqrt{0.002551} \approx \boxed{0.0505}
    \]
    Weights are sampled from Gaussian distribution $\mathcal{N}(0, 0.0505^2)$.
    
    \item \textbf{Xavier Initialization (for Sigmoid/Tanh)}:
    \[
    \text{limit} = \sqrt{\frac{6}{784 + 128}} = \sqrt{\frac{6}{912}} = \sqrt{0.006579} \approx \boxed{0.0811}
    \]
    Weights are sampled uniformly from interval $\mathcal{U}(-0.0811, 0.0811)$.
\end{enumerate}
\end{numexample}


% ============================================================
% CHAPTER 3: ACTIVATION FUNCTIONS
% ============================================================
\chapter{Activation Functions}

Activation functions introduce non-linearity into the network. Without them, any composition of linear layers is still linear and the network cannot learn complex patterns.

\section{ReLU (Rectified Linear Unit)}

\begin{definition}[ReLU]
\begin{equation}
    \text{ReLU}(z) = \max(0, z) = \begin{cases} z & \text{if } z > 0 \\ 0 & \text{if } z \leq 0 \end{cases}
\end{equation}

Its derivative is:
\begin{equation}
    \frac{d}{dz}\text{ReLU}(z) = \begin{cases} 1 & \text{if } z > 0 \\ 0 & \text{if } z \leq 0 \end{cases}
\end{equation}
\end{definition}

\begin{numexample}
Given $\mathbf{z} = \begin{bmatrix} 1.5 & -0.8 & 3.4 & -2.1 \end{bmatrix}$:
\[
\text{ReLU}(\mathbf{z}) = \begin{bmatrix} 1.5 & 0 & 3.4 & 0 \end{bmatrix}
\]
The derivative:
\[
\text{ReLU}'(\mathbf{z}) = \begin{bmatrix} 1 & 0 & 1 & 0 \end{bmatrix}
\]
\end{numexample}

\begin{tikzpicture}
\begin{axis}[
    width=8cm, height=5cm,
    xlabel={$z$}, ylabel={$\text{ReLU}(z)$},
    grid=major, axis lines=middle,
    xmin=-4, xmax=4, ymin=-1, ymax=4,
    title={ReLU Activation}
]
\addplot[blue, thick, domain=-4:0] {0};
\addplot[blue, thick, domain=0:4] {x};
\end{axis}
\end{tikzpicture}

\section{Sigmoid}

\begin{definition}[Sigmoid]
\begin{equation}
    \sigma(z) = \frac{1}{1 + e^{-z}}
\end{equation}

Its derivative has an elegant form:
\begin{equation}
    \sigma'(z) = \sigma(z) \cdot (1 - \sigma(z))
\end{equation}
\end{definition}

\begin{proof}
Let $\sigma(z) = (1+e^{-z})^{-1}$. By the chain rule:
\begin{align*}
    \sigma'(z) &= -1 \cdot (1+e^{-z})^{-2} \cdot (-e^{-z}) \\
    &= \frac{e^{-z}}{(1+e^{-z})^2} \\
    &= \frac{1}{1+e^{-z}} \cdot \frac{e^{-z}}{1+e^{-z}} \\
    &= \sigma(z) \cdot \frac{(1+e^{-z}) - 1}{1+e^{-z}} \\
    &= \sigma(z) \cdot \bigl(1 - \sigma(z)\bigr) \qedhere
\end{align*}
\end{proof}

\begin{numexample}
For $z = 1.0$:
\begin{align*}
    \sigma(1.0) &= \frac{1}{1 + e^{-1.0}} = \frac{1}{1 + 0.3679} = \frac{1}{1.3679} = \boxed{0.7311} \\[6pt]
    \sigma'(1.0) &= 0.7311 \times (1 - 0.7311) = 0.7311 \times 0.2689 = \boxed{0.1966}
\end{align*}
\end{numexample}

\section{Tanh (Hyperbolic Tangent)}

\begin{definition}[Tanh]
\begin{equation}
    \tanh(z) = \frac{e^z - e^{-z}}{e^z + e^{-z}}
\end{equation}

Its derivative:
\begin{equation}
    \tanh'(z) = 1 - \tanh^2(z)
\end{equation}
\end{definition}

\begin{proof}
Using the quotient rule with $u = e^z - e^{-z}$ and $v = e^z + e^{-z}$:
\begin{align*}
    \tanh'(z) &= \frac{u'v - uv'}{v^2} = \frac{(e^z+e^{-z})(e^z+e^{-z}) - (e^z-e^{-z})(e^z-e^{-z})}{(e^z+e^{-z})^2} \\
    &= \frac{(e^z+e^{-z})^2 - (e^z-e^{-z})^2}{(e^z+e^{-z})^2} \\
    &= 1 - \left(\frac{e^z - e^{-z}}{e^z + e^{-z}}\right)^2 = 1 - \tanh^2(z) \qedhere
\end{align*}
\end{proof}

\begin{numexample}
For $z = 0.5$:
\begin{align*}
    \tanh(0.5) &= \frac{e^{0.5} - e^{-0.5}}{e^{0.5} + e^{-0.5}} = \frac{1.6487 - 0.6065}{1.6487 + 0.6065} = \frac{1.0422}{2.2553} = \boxed{0.4621} \\[6pt]
    \tanh'(0.5) &= 1 - (0.4621)^2 = 1 - 0.2135 = \boxed{0.7864}
\end{align*}
\end{numexample}

\section{Softmax}

\begin{definition}[Softmax]
For a vector $\mathbf{z} \in \mathbb{R}^K$, the Softmax function outputs a probability distribution:
\begin{equation}
    \text{Softmax}(\mathbf{z})_i = \frac{e^{z_i}}{\sum_{j=1}^{K} e^{z_j}}, \quad i = 1, \ldots, K
\end{equation}

Properties: $\text{Softmax}(\mathbf{z})_i \in (0, 1)$ and $\sum_i \text{Softmax}(\mathbf{z})_i = 1$.
\end{definition}

The Jacobian matrix of Softmax is:
\begin{equation}
    \frac{\partial \text{Softmax}(\mathbf{z})_i}{\partial z_j} = \begin{cases}
        s_i(1 - s_i) & \text{if } i = j \\
        -s_i \cdot s_j & \text{if } i \neq j
    \end{cases}
\end{equation}
where $s_i = \text{Softmax}(\mathbf{z})_i$. In matrix form:
\begin{equation}
    J = \text{diag}(\mathbf{s}) - \mathbf{s}\mathbf{s}^\top
\end{equation}

\begin{numexample}
For $\mathbf{z} = \begin{bmatrix} 2.0 & 1.0 & 0.1 \end{bmatrix}$:
\begin{align*}
    e^{z_1} &= e^{2.0} = 7.389, \quad e^{z_2} = e^{1.0} = 2.718, \quad e^{z_3} = e^{0.1} = 1.105 \\
    \sum &= 7.389 + 2.718 + 1.105 = 11.212
\end{align*}
\[
\text{Softmax}(\mathbf{z}) = \begin{bmatrix} \frac{7.389}{11.212} & \frac{2.718}{11.212} & \frac{1.105}{11.212} \end{bmatrix} = \boxed{\begin{bmatrix} 0.659 & 0.242 & 0.099 \end{bmatrix}}
\]
Notice: $0.659 + 0.242 + 0.099 = 1.0$ \checkmark
\end{numexample}

\begin{keyidea}
In practice, we subtract $\max(\mathbf{z})$ before computing Softmax for numerical stability:
\[
\text{Softmax}(\mathbf{z})_i = \frac{e^{z_i - \max(\mathbf{z})}}{\sum_j e^{z_j - \max(\mathbf{z})}}
\]
This prevents overflow when $z_i$ is very large.
\end{keyidea}


% ============================================================
% CHAPTER 4: LOSS FUNCTIONS
% ============================================================
\chapter{Loss Functions}

\section{Categorical Cross-Entropy}

\begin{definition}[Categorical Cross-Entropy]
For a single sample with predicted probability vector $\hat{\mathbf{y}} \in \mathbb{R}^K$ and true one-hot label $\mathbf{y} \in \{0,1\}^K$:
\begin{equation}
    \mathcal{L}(\hat{\mathbf{y}}, \mathbf{y}) = -\sum_{i=1}^{K} y_i \ln(\hat{y}_i)
\end{equation}

Since $\mathbf{y}$ is one-hot (only one $y_c = 1$, all others $0$), this simplifies to:
\begin{equation}
    \mathcal{L} = -\ln(\hat{y}_c)
\end{equation}
where $c$ is the index of the true class.
\end{definition}

\begin{numexample}
Suppose the true class is $c = 0$ (the first class), and Softmax output is:
\[
\hat{\mathbf{y}} = \begin{bmatrix} 0.659 & 0.242 & 0.099 \end{bmatrix}
\]
One-hot label: $\mathbf{y} = \begin{bmatrix} 1 & 0 & 0 \end{bmatrix}$
\[
\mathcal{L} = -\ln(0.659) = \boxed{0.417}
\]
If the model were more confident ($\hat{y}_0 = 0.95$): $\mathcal{L} = -\ln(0.95) = 0.051$ (lower loss).\\
If the model were wrong ($\hat{y}_0 = 0.01$): $\mathcal{L} = -\ln(0.01) = 4.605$ (high loss).
\end{numexample}

\section{Combined Softmax + Cross-Entropy Gradient}

\begin{theorem}[Gradient of Cross-Entropy w.r.t.\ Softmax Input]
When using Softmax activation with Categorical Cross-Entropy loss, the gradient of the loss with respect to the pre-activation logits $\mathbf{z}$ is remarkably simple:
\begin{equation}
    \frac{\partial \mathcal{L}}{\partial z_i} = \hat{y}_i - y_i
\end{equation}
where $\hat{y}_i = \text{Softmax}(\mathbf{z})_i$.
\end{theorem}

\begin{proof}
We need $\frac{\partial \mathcal{L}}{\partial z_i}$. By chain rule:
\[
\frac{\partial \mathcal{L}}{\partial z_i} = \sum_{j=1}^{K} \frac{\partial \mathcal{L}}{\partial \hat{y}_j} \cdot \frac{\partial \hat{y}_j}{\partial z_i}
\]

The cross-entropy gradient w.r.t.\ $\hat{y}_j$ is:
\[
\frac{\partial \mathcal{L}}{\partial \hat{y}_j} = -\frac{y_j}{\hat{y}_j}
\]

The Softmax Jacobian gives:
\[
\frac{\partial \hat{y}_j}{\partial z_i} = \hat{y}_j(\delta_{ij} - \hat{y}_i)
\]
where $\delta_{ij}$ is the Kronecker delta.

Combining:
\begin{align*}
\frac{\partial \mathcal{L}}{\partial z_i} &= -\sum_{j} \frac{y_j}{\hat{y}_j} \cdot \hat{y}_j(\delta_{ij} - \hat{y}_i) \\
&= -\sum_{j} y_j (\delta_{ij} - \hat{y}_i) \\
&= -y_i + \hat{y}_i \sum_{j} y_j \\
&= -y_i + \hat{y}_i \cdot 1 \quad (\text{since } \sum_j y_j = 1) \\
&= \hat{y}_i - y_i \qedhere
\end{align*}
\end{proof}

\begin{numexample}
Continuing with $\hat{\mathbf{y}} = \begin{bmatrix} 0.659 & 0.242 & 0.099 \end{bmatrix}$ and $\mathbf{y} = \begin{bmatrix} 1 & 0 & 0 \end{bmatrix}$:
\[
\frac{\partial \mathcal{L}}{\partial \mathbf{z}} = \hat{\mathbf{y}} - \mathbf{y} = \begin{bmatrix} 0.659 - 1 \\ 0.242 - 0 \\ 0.099 - 0 \end{bmatrix} = \boxed{\begin{bmatrix} -0.341 \\ 0.242 \\ 0.099 \end{bmatrix}}
\]
Notice: the gradient for the correct class is negative (pushing it toward 1), while gradients for wrong classes are positive (pushing them toward 0).
\end{numexample}


% ============================================================
% CHAPTER 5: BACKPROPAGATION
% ============================================================
\chapter{Backpropagation}

\section{The Chain Rule}

Backpropagation is simply the \textbf{chain rule of calculus} applied systematically from the output layer back to the input layer.

\begin{theorem}[Multivariate Chain Rule]
If $\mathcal{L}$ depends on $\mathbf{z}$ through an intermediate variable $\mathbf{a} = \sigma(\mathbf{z})$, then:
\begin{equation}
    \frac{\partial \mathcal{L}}{\partial z_i} = \sum_j \frac{\partial \mathcal{L}}{\partial a_j} \cdot \frac{\partial a_j}{\partial z_i}
\end{equation}
\end{theorem}

\section{Backpropagation Through a Dense Layer}

Consider a Dense layer: $\mathbf{z} = X W + \mathbf{b}$.

Given the upstream gradient $\frac{\partial \mathcal{L}}{\partial \mathbf{z}} \in \mathbb{R}^{N \times d_\text{out}}$ (which we call \texttt{dvalues}), we need to compute:

\begin{enumerate}[label=\textbf{\arabic*.}]
    \item \textbf{Gradient w.r.t.\ weights} $\frac{\partial \mathcal{L}}{\partial W}$:
    \begin{equation}
        \frac{\partial \mathcal{L}}{\partial W} = X^\top \cdot \frac{\partial \mathcal{L}}{\partial \mathbf{z}}
    \end{equation}
    
    \item \textbf{Gradient w.r.t.\ biases} $\frac{\partial \mathcal{L}}{\partial \mathbf{b}}$:
    \begin{equation}
        \frac{\partial \mathcal{L}}{\partial \mathbf{b}} = \sum_{n=1}^{N} \frac{\partial \mathcal{L}}{\partial z_n} \quad \text{(sum over batch dimension)}
    \end{equation}
    
    \item \textbf{Gradient w.r.t.\ inputs} $\frac{\partial \mathcal{L}}{\partial X}$ (to pass to previous layer):
    \begin{equation}
        \frac{\partial \mathcal{L}}{\partial X} = \frac{\partial \mathcal{L}}{\partial \mathbf{z}} \cdot W^\top
    \end{equation}
\end{enumerate}

\begin{proof}[Rigorously Detailed Proof of Dense Layer Matrix Derivatives]
Let $X \in \mathbb{R}^{N \times d_\text{in}}$, $W \in \mathbb{R}^{d_\text{in} \times d_\text{out}}$, $\mathbf{b} \in \mathbb{R}^{d_\text{out}}$, and $Z = X W + \mathbf{b} \in \mathbb{R}^{N \times d_\text{out}}$.

The element-wise equation for entry $(i, j)$ in matrix $Z$ is:
\begin{equation}
    Z_{ij} = \sum_{k=1}^{d_\text{in}} X_{ik} W_{kj} + b_j
\end{equation}

\textbf{1. Weight Gradient Matrix Proof} ($\frac{\partial \mathcal{L}}{\partial W}$):
By the multivariate chain rule, the gradient of scalar loss $\mathcal{L}$ with respect to weight entry $W_{kj}$ is:
\begin{equation}
    \frac{\partial \mathcal{L}}{\partial W_{kj}} = \sum_{i=1}^{N} \sum_{l=1}^{d_\text{out}} \frac{\partial \mathcal{L}}{\partial Z_{il}} \frac{\partial Z_{il}}{\partial W_{kj}}
\end{equation}
Notice that $\frac{\partial Z_{il}}{\partial W_{kj}} \neq 0$ ONLY when $l = j$. Specifically:
\begin{equation}
    \frac{\partial Z_{ij}}{\partial W_{kj}} = X_{ik}
\end{equation}
Substituting this back into the sum simplifies it to:
\begin{equation}
    \frac{\partial \mathcal{L}}{\partial W_{kj}} = \sum_{i=1}^{N} \frac{\partial \mathcal{L}}{\partial Z_{ij}} X_{ik} = \sum_{i=1}^{N} X_{ik} \frac{\partial \mathcal{L}}{\partial Z_{ij}} = \sum_{i=1}^{N} (X^\top)_{ki} \left(\frac{\partial \mathcal{L}}{\partial Z}\right)_{ij}
\end{equation}
By the definition of matrix multiplication, entry $(k, j)$ of $(X^\top \cdot \frac{\partial \mathcal{L}}{\partial Z})$ is precisely $\sum_{i=1}^N (X^\top)_{ki} (\frac{\partial \mathcal{L}}{\partial Z})_{ij}$. Thus:
\begin{equation}
    \boxed{\frac{\partial \mathcal{L}}{\partial W} = X^\top \cdot \frac{\partial \mathcal{L}}{\partial Z}}
\end{equation}

\textbf{2. Input Gradient Matrix Proof} ($\frac{\partial \mathcal{L}}{\partial X}$):
Similarly, for input entry $X_{ik}$:
\begin{equation}
    \frac{\partial \mathcal{L}}{\partial X_{ik}} = \sum_{j=1}^{d_\text{out}} \frac{\partial \mathcal{L}}{\partial Z_{ij}} \frac{\partial Z_{ij}}{\partial X_{ik}}
\end{equation}
Since $\frac{\partial Z_{ij}}{\partial X_{ik}} = W_{kj}$:
\begin{equation}
    \frac{\partial \mathcal{L}}{\partial X_{ik}} = \sum_{j=1}^{d_\text{out}} \frac{\partial \mathcal{L}}{\partial Z_{ij}} W_{kj} = \sum_{j=1}^{d_\text{out}} \left(\frac{\partial \mathcal{L}}{\partial Z}\right)_{ij} (W^\top)_{jk}
\end{equation}
This matches entry $(i, k)$ of $(\frac{\partial \mathcal{L}}{\partial Z} \cdot W^\top)$. Therefore:
\begin{equation}
    \boxed{\frac{\partial \mathcal{L}}{\partial X} = \frac{\partial \mathcal{L}}{\partial Z} \cdot W^\top}
\end{equation}

\textbf{3. Bias Gradient Vector Proof} ($\frac{\partial \mathcal{L}}{\partial \mathbf{b}}$):
Since $\frac{\partial Z_{ij}}{\partial b_j} = 1$:
\begin{equation}
    \frac{\partial \mathcal{L}}{\partial b_j} = \sum_{i=1}^{N} \frac{\partial \mathcal{L}}{\partial Z_{ij}} \cdot 1 \implies \boxed{\frac{\partial \mathcal{L}}{\partial \mathbf{b}} = \sum_{i=1}^{N} \frac{\partial \mathcal{L}}{\partial Z_{i,:}}} \qedhere
\end{equation}
\end{proof}

\section{Complete Backpropagation Example}

Let us trace forward and backward through a \textbf{complete 2-layer network} with concrete numbers.

\subsection{Network Setup}

\begin{itemize}
    \item Input: $\mathbf{x} = \begin{bmatrix} 0.5 & 0.3 \end{bmatrix}$ (2 features)
    \item Hidden layer: 2 neurons with ReLU
    \item Output layer: 2 neurons with Softmax
    \item True label: class 0, so $\mathbf{y} = \begin{bmatrix} 1 & 0 \end{bmatrix}$
\end{itemize}

Parameters:
\begin{align*}
W^{[1]} &= \begin{bmatrix} 0.15 & 0.25 \\ 0.20 & 0.30 \end{bmatrix}, \quad
\mathbf{b}^{[1]} = \begin{bmatrix} 0.35 & 0.35 \end{bmatrix} \\
W^{[2]} &= \begin{bmatrix} 0.40 & 0.50 \\ 0.45 & 0.55 \end{bmatrix}, \quad
\mathbf{b}^{[2]} = \begin{bmatrix} 0.60 & 0.60 \end{bmatrix}
\end{align*}

\begin{figure}[H]
\centering
\begin{tikzpicture}[
    node distance=2.8cm and 3.5cm,
    input_node/.style={circle, draw=blue!80, fill=blue!10, thick, minimum size=1.1cm, align=center},
    hidden_node/.style={circle, draw=orange!80, fill=orange!10, thick, minimum size=1.3cm, align=center},
    output_node/.style={circle, draw=green!80, fill=green!10, thick, minimum size=1.3cm, align=center},
    forward_arr/.style={->, >=stealth, thick, blue!70!black},
    backward_arr/.style={<-, >=stealth, dashed, thick, red!70!black}
]

    % Nodes - Input Layer
    \node[input_node] (x1) at (0, 1.8) {$x_1$\\ \tiny $0.5$};
    \node[input_node] (x2) at (0, -1.8) {$x_2$\\ \tiny $0.3$};

    % Nodes - Hidden Layer (ReLU)
    \node[hidden_node] (h1) at (4, 1.8) {$a_1^{[1]}$\\ \tiny $0.485$};
    \node[hidden_node] (h2) at (4, -1.8) {$a_2^{[1]}$\\ \tiny $0.565$};

    % Nodes - Output Layer (Softmax)
    \node[output_node] (o1) at (8, 1.8) {$\hat{y}_1$\\ \tiny $0.474$};
    \node[output_node] (o2) at (8, -1.8) {$\hat{y}_2$\\ \tiny $0.526$};

    % Forward Connections Layer 1
    \draw[forward_arr] (x1) -- node[above, sloped, font=\tiny] {$w_{11}^{[1]}=0.15$} (h1);
    \draw[forward_arr] (x1) -- node[near start, above, sloped, font=\tiny] {$w_{12}^{[1]}=0.25$} (h2);
    \draw[forward_arr] (x2) -- node[near start, below, sloped, font=\tiny] {$w_{21}^{[1]}=0.20$} (h1);
    \draw[forward_arr] (x2) -- node[below, sloped, font=\tiny] {$w_{22}^{[1]}=0.30$} (h2);

    % Forward Connections Layer 2
    \draw[forward_arr] (h1) -- node[above, sloped, font=\tiny] {$w_{11}^{[2]}=0.40$} (o1);
    \draw[forward_arr] (h1) -- node[near start, above, sloped, font=\tiny] {$w_{12}^{[2]}=0.50$} (o2);
    \draw[forward_arr] (h2) -- node[near start, below, sloped, font=\tiny] {$w_{21}^{[2]}=0.45$} (o1);
    \draw[forward_arr] (h2) -- node[below, sloped, font=\tiny] {$w_{22}^{[2]}=0.55$} (o2);

    % Backward Gradient Flow Indicators (Dashed Red Lines)
    \draw[backward_arr] ($(h1)+(0,0.8)$) -- node[above, font=\tiny, color=red!80!black] {$\frac{\partial \mathcal{L}}{\partial \mathbf{a}^{[1]}} = 0.053$} ($(o1)+(0,0.8)$);
    \draw[backward_arr] ($(x1)+(0,0.8)$) -- node[above, font=\tiny, color=red!80!black] {$\frac{\partial \mathcal{L}}{\partial \mathbf{z}^{[1]}} = 0.053$} ($(h1)+(0,0.8)$);

    % Biases
    \node[above=0.1cm of h1, font=\tiny, color=orange!90!black] {$b_1^{[1]}=0.35$};
    \node[below=0.1cm of h2, font=\tiny, color=orange!90!black] {$b_2^{[1]}=0.35$};
    \node[above=0.1cm of o1, font=\tiny, color=green!90!black] {$b_1^{[2]}=0.60$};
    \node[below=0.1cm of o2, font=\tiny, color=green!90!black] {$b_2^{[2]}=0.60$};

    % Layer Headers
    \node[above=1.2cm of x1, font=\bfseries\small] {Input Layer};
    \node[above=1.2cm of h1, font=\bfseries\small] {Hidden Layer (ReLU)};
    \node[above=1.2cm of o1, font=\bfseries\small] {Output Layer (Softmax)};

\end{tikzpicture}
\caption{Neural network diagram for the 2-layer backpropagation example. Blue arrows indicate forward propagation values; red dashed arrows indicate the backward flow of loss gradients.}
\label{fig:backprop_nn_diagram}
\end{figure}

\subsection{Forward Pass}

\textbf{Layer 1 (Dense + ReLU):}
\begin{align*}
z_1^{[1]} &= (0.5)(0.15) + (0.3)(0.20) + 0.35 = 0.075 + 0.060 + 0.35 = 0.485 \\
z_2^{[1]} &= (0.5)(0.25) + (0.3)(0.30) + 0.35 = 0.125 + 0.090 + 0.35 = 0.565
\end{align*}
\[
\mathbf{z}^{[1]} = \begin{bmatrix} 0.485 & 0.565 \end{bmatrix}
\]

Apply ReLU (both values are positive, so unchanged):
\[
\mathbf{a}^{[1]} = \text{ReLU}(\mathbf{z}^{[1]}) = \begin{bmatrix} 0.485 & 0.565 \end{bmatrix}
\]

\textbf{Layer 2 (Dense + Softmax):}
\begin{align*}
z_1^{[2]} &= (0.485)(0.40) + (0.565)(0.45) + 0.60 = 0.194 + 0.254 + 0.60 = 1.048 \\
z_2^{[2]} &= (0.485)(0.50) + (0.565)(0.55) + 0.60 = 0.2425 + 0.3108 + 0.60 = 1.153
\end{align*}
\[
\mathbf{z}^{[2]} = \begin{bmatrix} 1.048 & 1.153 \end{bmatrix}
\]

Apply Softmax:
\begin{align*}
e^{1.048} &= 2.852, \quad e^{1.153} = 3.168 \\
\sum &= 2.852 + 3.168 = 6.020
\end{align*}
\[
\hat{\mathbf{y}} = \mathbf{a}^{[2]} = \begin{bmatrix} \frac{2.852}{6.020} & \frac{3.168}{6.020} \end{bmatrix} = \begin{bmatrix} 0.474 & 0.526 \end{bmatrix}
\]

\textbf{Loss:}
\[
\mathcal{L} = -\ln(\hat{y}_0) = -\ln(0.474) = \boxed{0.747}
\]

\subsection{Backward Pass}

\textbf{Step 1: Gradient at Softmax + Cross-Entropy output} (using the combined formula):
\[
\frac{\partial \mathcal{L}}{\partial \mathbf{z}^{[2]}} = \hat{\mathbf{y}} - \mathbf{y} = \begin{bmatrix} 0.474 - 1 & 0.526 - 0 \end{bmatrix} = \begin{bmatrix} -0.526 & 0.526 \end{bmatrix}
\]

\textbf{Step 2: Backprop through Dense Layer 2:}

Weight gradients:
\[
\frac{\partial \mathcal{L}}{\partial W^{[2]}} = (\mathbf{a}^{[1]})^\top \cdot \frac{\partial \mathcal{L}}{\partial \mathbf{z}^{[2]}} = \begin{bmatrix} 0.485 \\ 0.565 \end{bmatrix} \begin{bmatrix} -0.526 & 0.526 \end{bmatrix} = \begin{bmatrix} -0.255 & 0.255 \\ -0.297 & 0.297 \end{bmatrix}
\]

Bias gradients:
\[
\frac{\partial \mathcal{L}}{\partial \mathbf{b}^{[2]}} = \begin{bmatrix} -0.526 & 0.526 \end{bmatrix}
\]

Input gradients (to pass to Layer 1):
\[
\frac{\partial \mathcal{L}}{\partial \mathbf{a}^{[1]}} = \frac{\partial \mathcal{L}}{\partial \mathbf{z}^{[2]}} \cdot (W^{[2]})^\top = \begin{bmatrix} -0.526 & 0.526 \end{bmatrix} \begin{bmatrix} 0.40 & 0.45 \\ 0.50 & 0.55 \end{bmatrix} = \begin{bmatrix} 0.053 & 0.053 \end{bmatrix}
\]
Computing element by element:
\begin{align*}
    \delta_1 &= (-0.526)(0.40) + (0.526)(0.50) = -0.210 + 0.263 = 0.053 \\
    \delta_2 &= (-0.526)(0.45) + (0.526)(0.55) = -0.237 + 0.289 = 0.053
\end{align*}

\textbf{Step 3: Backprop through ReLU:}

Since both $z_1^{[1]} = 0.485 > 0$ and $z_2^{[1]} = 0.565 > 0$:
\[
\frac{\partial \mathcal{L}}{\partial \mathbf{z}^{[1]}} = \frac{\partial \mathcal{L}}{\partial \mathbf{a}^{[1]}} \odot \text{ReLU}'(\mathbf{z}^{[1]}) = \begin{bmatrix} 0.053 & 0.053 \end{bmatrix} \odot \begin{bmatrix} 1 & 1 \end{bmatrix} = \begin{bmatrix} 0.053 & 0.053 \end{bmatrix}
\]

\textbf{Step 4: Backprop through Dense Layer 1:}

Weight gradients:
\[
\frac{\partial \mathcal{L}}{\partial W^{[1]}} = \mathbf{x}^\top \cdot \frac{\partial \mathcal{L}}{\partial \mathbf{z}^{[1]}} = \begin{bmatrix} 0.5 \\ 0.3 \end{bmatrix} \begin{bmatrix} 0.053 & 0.053 \end{bmatrix} = \begin{bmatrix} 0.026 & 0.026 \\ 0.016 & 0.016 \end{bmatrix}
\]

Bias gradients:
\[
\frac{\partial \mathcal{L}}{\partial \mathbf{b}^{[1]}} = \begin{bmatrix} 0.053 & 0.053 \end{bmatrix}
\]

\section{Binary Classification Backpropagation (Sigmoid + BCE)}

For binary classification tasks, the output layer uses a single neuron with Sigmoid activation $\hat{y} = \sigma(z) = \frac{1}{1+e^{-z}}$ and Binary Cross-Entropy (BCE) loss:
\begin{equation}
    \mathcal{L}(y, \hat{y}) = - \bigl[ y \ln(\hat{y}) + (1 - y) \ln(1 - \hat{y}) \bigr]
\end{equation}

\begin{theorem}[Gradient of BCE w.r.t.\ Pre-Activation Logit $z$]
The derivative of Binary Cross-Entropy loss with respect to logit $z$ is:
\begin{equation}
    \frac{\partial \mathcal{L}}{\partial z} = \hat{y} - y
\end{equation}
\end{theorem}

\begin{proof}
Applying the chain rule:
\[
\frac{\partial \mathcal{L}}{\partial z} = \frac{\partial \mathcal{L}}{\partial \hat{y}} \cdot \frac{\partial \hat{y}}{\partial z}
\]
1. Derivative of loss w.r.t.\ prediction $\hat{y}$:
\[
\frac{\partial \mathcal{L}}{\partial \hat{y}} = -\frac{y}{\hat{y}} + \frac{1 - y}{1 - \hat{y}} = \frac{-\hat{y}y(1-\hat{y}) + (1-y)\hat{y}}{\hat{y}(1-\hat{y})} = \frac{\hat{y} - y}{\hat{y}(1-\hat{y})}
\]
2. Derivative of Sigmoid w.r.t.\ logit $z$:
\[
\frac{\partial \hat{y}}{\partial z} = \sigma'(z) = \hat{y}(1 - \hat{y})
\]
3. Combining them via chain rule:
\[
\frac{\partial \mathcal{L}}{\partial z} = \frac{\hat{y} - y}{\hat{y}(1-\hat{y})} \cdot \hat{y}(1 - \hat{y}) = \boxed{\hat{y} - y} \qedhere
\]
\end{proof}

\begin{numexample}[Binary Classification Backpropagation]
Consider a single-neuron model for binary classification:
Input $\mathbf{x} = \begin{bmatrix} 1.0 & 2.0 \end{bmatrix}$, Weights $\mathbf{w} = \begin{bmatrix} 0.5 & -0.5 \end{bmatrix}$, Bias $b = 0.2$, True target $y = 1.0$.

\textbf{Forward Pass}:
\begin{align*}
z &= w_1 x_1 + w_2 x_2 + b = (0.5)(1.0) + (-0.5)(2.0) + 0.2 = 0.5 - 1.0 + 0.2 = -0.3 \\
\hat{y} &= \sigma(-0.3) = \frac{1}{1 + e^{0.3}} = \frac{1}{1 + 1.3499} = \boxed{0.4256}
\end{align*}

\textbf{Binary Cross-Entropy Loss}:
\[
\mathcal{L} = -\ln(\hat{y}) = -\ln(0.4256) = \boxed{0.8542}
\]

\textbf{Backward Pass}:
\begin{align*}
\frac{\partial \mathcal{L}}{\partial z} &= \hat{y} - y = 0.4256 - 1.0 = \boxed{-0.5744} \\[4pt]
\frac{\partial \mathcal{L}}{\partial \mathbf{w}} &= \mathbf{x}^\top \cdot \frac{\partial \mathcal{L}}{\partial z} = \begin{bmatrix} 1.0 \\ 2.0 \end{bmatrix} (-0.5744) = \boxed{\begin{bmatrix} -0.5744 \\ -1.1488 \end{bmatrix}} \\[4pt]
\frac{\partial \mathcal{L}}{\partial b} &= \frac{\partial \mathcal{L}}{\partial z} = \boxed{-0.5744}
\end{align*}

\textbf{Parameter Update ($\eta = 0.1$)}:
\begin{align*}
\mathbf{w}_\text{new} &= \begin{bmatrix} 0.5 & -0.5 \end{bmatrix} - 0.1 \cdot \begin{bmatrix} -0.5744 & -1.1488 \end{bmatrix} = \boxed{\begin{bmatrix} 0.5574 & -0.3851 \end{bmatrix}} \\
b_\text{new} &= 0.2 - 0.1 \cdot (-0.5744) = \boxed{0.2574}
\end{align*}
\end{numexample}

\section{Scalar Autograd Computational Graph Example}

In `mini-keras/data.py` (inspired by Andrej Karpathy's `micrograd`), backpropagation operates on individual scalar `Data` objects. Let us trace a complete scalar computation graph:

Consider inputs $x_1 = 2.0$, $x_2 = 0.0$, weights $w_1 = -3.0$, $w_2 = 1.0$, and bias $b = 6.88137$.

\textbf{Step 1: Forward Graph Computation}:
\begin{align*}
    u_1 &= x_1 \cdot w_1 = (2.0)(-3.0) = -6.0 \\
    u_2 &= x_2 \cdot w_2 = (0.0)(1.0) = 0.0 \\
    u_3 &= u_1 + u_2 = -6.0 + 0.0 = -6.0 \\
    y &= u_3 + b = -6.0 + 6.88137 = \boxed{0.88137}
\end{align*}

\textbf{Step 2: Backward Pass (Analytical Derivatives)}:
Start at root node: $\frac{\partial y}{\partial y} = 1.0$.

1. Addition node $y = u_3 + b$:
   \[
   \frac{\partial y}{\partial u_3} = 1.0 \cdot 1.0 = \boxed{1.0}, \quad \frac{\partial y}{\partial b} = 1.0 \cdot 1.0 = \boxed{1.0}
   \]
2. Addition node $u_3 = u_1 + u_2$:
   \[
   \frac{\partial y}{\partial u_1} = \frac{\partial y}{\partial u_3} \cdot 1 = \boxed{1.0}, \quad \frac{\partial y}{\partial u_2} = \frac{\partial y}{\partial u_3} \cdot 1 = \boxed{1.0}
   \]
3. Multiplication node $u_1 = x_1 \cdot w_1$:
   \[
   \frac{\partial y}{\partial x_1} = \frac{\partial y}{\partial u_1} \cdot w_1 = (1.0)(-3.0) = \boxed{-3.0}, \quad \frac{\partial y}{\partial w_1} = \frac{\partial y}{\partial u_1} \cdot x_1 = (1.0)(2.0) = \boxed{2.0}
   \]
4. Multiplication node $u_2 = x_2 \cdot w_2$:
   \[
   \frac{\partial y}{\partial x_2} = \frac{\partial y}{\partial u_2} \cdot w_2 = (1.0)(1.0) = \boxed{1.0}, \quad \frac{\partial y}{\partial w_2} = \frac{\partial y}{\partial u_2} \cdot x_2 = (1.0)(0.0) = \boxed{0.0}
   \]

\textbf{Summary Table of Scalar Autograd Gradients}:
\begin{center}
\renewcommand{\arraystretch}{1.3}
\begin{tabular}{cccc}
\toprule
\textbf{Variable} & \textbf{Forward Value} & \textbf{Analytical Gradient $\frac{\partial y}{\partial \text{var}}$} & \textbf{Meaning} \\
\midrule
$x_1$ & $2.0$ & $-3.0$ & Increasing $x_1$ by $1$ decreases $y$ by $3$ \\
$w_1$ & $-3.0$ & $2.0$ & Increasing $w_1$ by $1$ increases $y$ by $2$ \\
$x_2$ & $0.0$ & $1.0$ & Increasing $x_2$ by $1$ increases $y$ by $1$ \\
$w_2$ & $1.0$ & $0.0$ & $x_2=0$ so $w_2$ has no effect on $y$ \\
$b$ & $6.88137$ & $1.0$ & Increasing $b$ by $1$ increases $y$ by $1$ \\
\bottomrule
\end{tabular}
\end{center}

\begin{keyidea}
\textbf{Summary of backpropagation through each layer type:}

\begin{center}
\renewcommand{\arraystretch}{1.5}
\begin{tabular}{lcc}
\toprule
\textbf{Layer} & \textbf{Forward} & \textbf{Backward (given $\frac{\partial \mathcal{L}}{\partial \text{output}}$)} \\
\midrule
Dense & $Z = XW + \mathbf{b}$ & $\frac{\partial \mathcal{L}}{\partial W} = X^\top \frac{\partial \mathcal{L}}{\partial Z}$, \; $\frac{\partial \mathcal{L}}{\partial X} = \frac{\partial \mathcal{L}}{\partial Z} W^\top$ \\
ReLU & $A = \max(0, Z)$ & $\frac{\partial \mathcal{L}}{\partial Z} = \frac{\partial \mathcal{L}}{\partial A} \odot \mathbb{1}_{Z > 0}$ \\
Softmax+CE & $\hat{y} = \text{Softmax}(Z)$ & $\frac{\partial \mathcal{L}}{\partial Z} = \hat{\mathbf{y}} - \mathbf{y}$ \\
\bottomrule
\end{tabular}
\end{center}
\end{keyidea}


% ============================================================
% CHAPTER 6: OPTIMIZATION ALGORITHMS
% ============================================================
\chapter{Optimization Algorithms}

After computing gradients via backpropagation, we need an algorithm to \textit{update} the parameters. All optimizers follow the general pattern:
\begin{equation}
    \theta \leftarrow \theta - \eta \cdot \Delta\theta
\end{equation}
where $\eta$ is the learning rate and $\Delta\theta$ is the update direction.

\section{Stochastic Gradient Descent (SGD)}

The simplest optimizer. Update rule:
\begin{equation}
    \theta_{t+1} = \theta_t - \eta \cdot \nabla_\theta \mathcal{L}
\end{equation}

Expanded for weights and biases:
\begin{align}
    W &\leftarrow W - \eta \cdot \frac{\partial \mathcal{L}}{\partial W} \\
    \mathbf{b} &\leftarrow \mathbf{b} - \eta \cdot \frac{\partial \mathcal{L}}{\partial \mathbf{b}}
\end{align}

\begin{numexample}[Complete Parameter Updates with SGD]
Using learning rate $\eta = 0.5$ and the gradients calculated from Chapter 5:

\textbf{Layer 2 Updates}:
\begin{align*}
W^{[2]}_\text{new} &= \begin{bmatrix} 0.40 & 0.50 \\ 0.45 & 0.55 \end{bmatrix} - 0.5 \cdot \begin{bmatrix} -0.255 & 0.255 \\ -0.297 & 0.297 \end{bmatrix} = \boxed{\begin{bmatrix} 0.528 & 0.372 \\ 0.599 & 0.401 \end{bmatrix}} \\[6pt]
\mathbf{b}^{[2]}_\text{new} &= \begin{bmatrix} 0.60 & 0.60 \end{bmatrix} - 0.5 \cdot \begin{bmatrix} -0.526 & 0.526 \end{bmatrix} = \boxed{\begin{bmatrix} 0.863 & 0.337 \end{bmatrix}}
\end{align*}

\textbf{Layer 1 Updates}:
\begin{align*}
W^{[1]}_\text{new} &= \begin{bmatrix} 0.15 & 0.25 \\ 0.20 & 0.30 \end{bmatrix} - 0.5 \cdot \begin{bmatrix} 0.026 & 0.026 \\ 0.016 & 0.016 \end{bmatrix} = \boxed{\begin{bmatrix} 0.137 & 0.237 \\ 0.192 & 0.292 \end{bmatrix}} \\[6pt]
\mathbf{b}^{[1]}_\text{new} &= \begin{bmatrix} 0.35 & 0.35 \end{bmatrix} - 0.5 \cdot \begin{bmatrix} 0.053 & 0.053 \end{bmatrix} = \boxed{\begin{bmatrix} 0.324 & 0.324 \end{bmatrix}}
\end{align*}
\end{numexample}

\begin{numexample}[Iteration 2 Forward Pass: Proving Loss Drop]
Now let us run the forward pass on Iteration 2 using these updated parameters for input $\mathbf{x} = \begin{bmatrix} 0.5 & 0.3 \end{bmatrix}$:

\textbf{Layer 1 Forward}:
\begin{align*}
z_1^{[1]} &= (0.5)(0.137) + (0.3)(0.192) + 0.324 = 0.0685 + 0.0576 + 0.324 = 0.4501 \\
z_2^{[1]} &= (0.5)(0.237) + (0.3)(0.292) + 0.324 = 0.1185 + 0.0876 + 0.324 = 0.5301
\end{align*}
Applying ReLU: $\mathbf{a}^{[1]} = \begin{bmatrix} 0.4501 & 0.5301 \end{bmatrix}$

\textbf{Layer 2 Forward}:
\begin{align*}
z_1^{[2]} &= (0.4501)(0.528) + (0.5301)(0.599) + 0.863 = 0.2377 + 0.3175 + 0.863 = 1.4182 \\
z_2^{[2]} &= (0.4501)(0.372) + (0.5301)(0.401) + 0.337 = 0.1674 + 0.2126 + 0.337 = 0.7170
\end{align*}
Applying Softmax:
\begin{align*}
e^{1.4182} &= 4.1297, \quad e^{0.7170} = 2.0483, \quad \sum = 6.1780 \\
\hat{\mathbf{y}} &= \begin{bmatrix} \frac{4.1297}{6.1780} & \frac{2.0483}{6.1780} \end{bmatrix} = \begin{bmatrix} 0.6685 & 0.3315 \end{bmatrix}
\end{align*}

\textbf{New Loss}:
\[
\mathcal{L}_2 = -\ln(0.6685) = \boxed{0.4027}
\]

\textbf{Mathematical Proof of Learning}:
\[
\mathcal{L}_1 = 0.747 \longrightarrow \mathcal{L}_2 = 0.4027 \quad \text{(\textbf{Loss reduced by 46\% in a single update!})}
\]
The predicted probability for the correct class increased from $47.4\%$ to $66.85\%$.
\end{numexample}

\section{SGD with Momentum}

Momentum adds ``inertia'' to updates, accelerating convergence:
\begin{align}
    \mathbf{v}_t &= \beta \cdot \mathbf{v}_{t-1} - \eta \cdot \nabla_\theta \mathcal{L} \\
    \theta_{t+1} &= \theta_t + \mathbf{v}_t
\end{align}
where $\beta \in [0, 1)$ is the momentum coefficient (typically $\beta = 0.9$) and $\mathbf{v}_0 = \mathbf{0}$.

\begin{keyidea}
Momentum is analogous to a ball rolling downhill. The velocity $\mathbf{v}$ accumulates past gradients, causing the optimizer to move faster in consistent gradient directions and dampen oscillations.
\end{keyidea}

\begin{numexample}[2-Step Momentum Calculation]
Consider a single scalar parameter $\theta_0 = 1.0$ with momentum $\beta = 0.9$ and learning rate $\eta = 0.1$.
Suppose the gradients over two consecutive steps are $g_1 = 0.5$ and $g_2 = 0.4$.

\textbf{Step 1 ($t=1$)}:
\begin{align*}
v_1 &= \beta v_0 - \eta g_1 = (0.9)(0) - (0.1)(0.5) = \boxed{-0.05} \\
\theta_1 &= \theta_0 + v_1 = 1.0 + (-0.05) = \boxed{0.95}
\end{align*}

\textbf{Step 2 ($t=2$)}:
\begin{align*}
v_2 &= \beta v_1 - \eta g_2 = (0.9)(-0.05) - (0.1)(0.4) = -0.045 - 0.040 = \boxed{-0.085} \\
\theta_2 &= \theta_1 + v_2 = 0.95 + (-0.085) = \boxed{0.865}
\end{align*}

Notice how velocity accumulated: $v_2$ is larger than $v_1$ because the gradient kept pointing in the same direction!
\end{numexample}

\section{AdaGrad}

AdaGrad \textit{adapts} the learning rate per-parameter based on historical gradient magnitudes:
\begin{align}
    G_t &= G_{t-1} + (\nabla_\theta \mathcal{L})^2 \quad \text{(element-wise squaring)} \\
    \theta_{t+1} &= \theta_t - \frac{\eta}{\sqrt{G_t} + \epsilon} \cdot \nabla_\theta \mathcal{L}
\end{align}
where $\epsilon \approx 10^{-7}$ prevents division by zero.

\begin{warning}
AdaGrad's cache $G_t$ monotonically increases, causing the effective learning rate to shrink to zero over time. This is why RMSprop and Adam were developed.
\end{warning}

\section{RMSprop}

RMSprop fixes AdaGrad by using an \textit{exponentially decaying average} of squared gradients:
\begin{align}
    G_t &= \rho \cdot G_{t-1} + (1 - \rho) \cdot (\nabla_\theta \mathcal{L})^2 \\
    \theta_{t+1} &= \theta_t - \frac{\eta}{\sqrt{G_t} + \epsilon} \cdot \nabla_\theta \mathcal{L}
\end{align}
where $\rho = 0.9$ (decay rate).

\begin{numexample}
Suppose $\nabla_\theta \mathcal{L} = 0.5$, $G_0 = 0$, $\rho = 0.9$, $\eta = 0.01$, $\epsilon = 10^{-7}$:

\textbf{Step 1}: $G_1 = 0.9 \cdot 0 + 0.1 \cdot (0.5)^2 = 0.025$

Update: $\Delta\theta = \frac{0.01}{\sqrt{0.025} + 10^{-7}} \cdot 0.5 = \frac{0.01}{0.1581} \cdot 0.5 = \boxed{0.0316}$

\textbf{Step 2} (same gradient): $G_2 = 0.9 \cdot 0.025 + 0.1 \cdot 0.025 = 0.025$

The cache stabilizes, unlike AdaGrad where it would keep growing.
\end{numexample}

\section{Adam (Adaptive Moment Estimation)}

Adam combines momentum (first moment) with RMSprop (second moment):
\begin{align}
    \mathbf{m}_t &= \beta_1 \mathbf{m}_{t-1} + (1-\beta_1) \nabla_\theta \mathcal{L} \quad \text{(first moment estimate)} \\
    \mathbf{v}_t &= \beta_2 \mathbf{v}_{t-1} + (1-\beta_2) (\nabla_\theta \mathcal{L})^2 \quad \text{(second moment estimate)} \\
    \hat{\mathbf{m}}_t &= \frac{\mathbf{m}_t}{1 - \beta_1^t} \quad \text{(bias correction)} \\
    \hat{\mathbf{v}}_t &= \frac{\mathbf{v}_t}{1 - \beta_2^t} \quad \text{(bias correction)} \\
    \theta_{t+1} &= \theta_t - \frac{\eta}{\sqrt{\hat{\mathbf{v}}_t} + \epsilon} \cdot \hat{\mathbf{m}}_t
\end{align}
Default hyperparameters: $\beta_1 = 0.9$, $\beta_2 = 0.999$, $\epsilon = 10^{-7}$.

\begin{numexample}[2-Step Adam Calculation]
Let $\theta_0 = 1.0$, learning rate $\eta = 0.01$, $\beta_1 = 0.9$, $\beta_2 = 0.999$, $\epsilon = 10^{-7}$.
Suppose gradients are $g_1 = 0.5$ at $t=1$ and $g_2 = 0.4$ at $t=2$.

\textbf{Step 1 ($t=1$)}:
\begin{align*}
    m_1 &= 0.9(0) + 0.1(0.5) = 0.05, \quad v_1 = 0.999(0) + 0.001(0.5^2) = 0.00025 \\
    \hat{m}_1 &= \frac{0.05}{1 - 0.9^1} = 0.5, \quad \hat{v}_1 = \frac{0.00025}{1 - 0.999^1} = 0.25 \\
    \theta_1 &= \theta_0 - \frac{0.01}{\sqrt{0.25} + 10^{-7}} \cdot 0.5 = 1.0 - 0.01 = \boxed{0.99}
\end{align*}

\textbf{Step 2 ($t=2$)}:
\begin{align*}
    m_2 &= 0.9(0.05) + 0.1(0.4) = 0.045 + 0.040 = 0.085 \\
    v_2 &= 0.999(0.00025) + 0.001(0.4^2) = 0.00024975 + 0.00016 = 0.00040975 \\[4pt]
    \hat{m}_2 &= \frac{0.085}{1 - 0.9^2} = \frac{0.085}{0.19} = 0.44737 \\
    \hat{v}_2 &= \frac{0.00040975}{1 - 0.999^2} = \frac{0.00040975}{0.001999} = 0.20498 \\[4pt]
    \theta_2 &= \theta_1 - \frac{0.01}{\sqrt{0.20498} + 10^{-7}} \cdot 0.44737 = 0.99 - \frac{0.01}{0.45275} \cdot 0.44737 \\
    &= 0.99 - 0.00988 = \boxed{0.98012}
\end{align*}

Note how bias correction factors $1-\beta_1^t$ ($0.1$ at $t=1 \to 0.19$ at $t=2$) dynamically adjust as $t$ increases.
\end{numexample}

\begin{keyidea}
\textbf{Comparison of Optimizers:}
\begin{center}
\renewcommand{\arraystretch}{1.4}
\begin{tabular}{lcccc}
\toprule
\textbf{Optimizer} & \textbf{Momentum} & \textbf{Adaptive LR} & \textbf{Bias Correction} & \textbf{Typical Use} \\
\midrule
SGD & \xmark & \xmark & \xmark & Simple tasks \\
SGD+Momentum & \checkmark & \xmark & \xmark & CNNs \\
AdaGrad & \xmark & \checkmark & \xmark & Sparse features \\
RMSprop & \xmark & \checkmark & \xmark & RNNs \\
Adam & \checkmark & \checkmark & \checkmark & \textbf{Default choice} \\
\bottomrule
\end{tabular}
\end{center}
\end{keyidea}


% ============================================================
% CHAPTER 7: PUTTING IT ALL TOGETHER
% ============================================================
\chapter{Putting It All Together: Mini-Keras}

\section{Training Loop}

The complete training algorithm:

\begin{enumerate}[label=\textbf{Step \arabic*:}]
    \item \textbf{Forward pass}: Compute $\hat{\mathbf{y}} = f(\mathbf{x}; \theta)$ through all layers.
    \item \textbf{Compute loss}: $\mathcal{L} = -\sum_i y_i \ln \hat{y}_i$.
    \item \textbf{Backward pass}: Compute $\nabla_\theta \mathcal{L}$ via chain rule through each layer in reverse.
    \item \textbf{Update parameters}: $\theta \leftarrow \theta - \eta \cdot \nabla_\theta \mathcal{L}$ (or Adam, etc.).
    \item \textbf{Repeat} for all mini-batches and epochs.
\end{enumerate}

\section{Mini-Batch Training}

For a dataset of $N$ samples with batch size $B$:
\begin{itemize}
    \item Number of batches per epoch: $\lfloor N / B \rfloor$
    \item Each batch computes the average loss over $B$ samples
    \item Gradients are normalized by $B$ to keep learning rate consistent
\end{itemize}

The average loss over a mini-batch:
\begin{equation}
    \mathcal{L}_\text{batch} = \frac{1}{B} \sum_{i=1}^{B} \mathcal{L}(\hat{\mathbf{y}}_i, \mathbf{y}_i)
\end{equation}

\section{Mini-Keras Code Mapping}

The following table shows how each mathematical concept maps to Mini-Keras code:

\begin{center}
\renewcommand{\arraystretch}{1.5}
\begin{tabular}{lll}
\toprule
\textbf{Math} & \textbf{Mini-Keras Code} & \textbf{Module} \\
\midrule
$Z = XW + \mathbf{b}$ & \texttt{Dense(units=128)} & \texttt{mini\_keras.layers} \\
$\text{ReLU}(Z)$ & \texttt{ReLU()} & \texttt{mini\_keras.activations} \\
$\text{Softmax}(Z)$ & \texttt{Softmax()} & \texttt{mini\_keras.activations} \\
$\mathcal{L} = -\sum y_i \ln \hat{y}_i$ & \texttt{CategoricalCrossentropy()} & \texttt{mini\_keras.losses} \\
$\theta \leftarrow \theta - \eta \nabla \mathcal{L}$ & \texttt{SGD(lr=0.01)} & \texttt{mini\_keras.optimizers} \\
Adam update & \texttt{Adam(lr=0.01)} & \texttt{mini\_keras.optimizers} \\
\bottomrule
\end{tabular}
\end{center}

\section{Example: MNIST Digit Classification}

\begin{lstlisting}
from mini_keras import Sequential, Dense
from mini_keras.activations import ReLU, Softmax
from mini_keras.datasets import load_mnist
from mini_keras.utils import to_categorical

# Load and preprocess MNIST
(X_train, y_train), (X_test, y_test) = load_mnist()
X_train = X_train.reshape(-1, 784).astype('float32') / 255.0
Y_train = to_categorical(y_train, num_classes=10)

# Build model (784 -> 128 -> 64 -> 10)
model = Sequential()
model.add(Dense(units=128, input_shape=(784,)))  # Z = X @ W + b
model.add(ReLU())                                 # A = max(0, Z)
model.add(Dense(units=64))
model.add(ReLU())
model.add(Dense(units=10))
model.add(Softmax())                              # y_hat = softmax(Z)

# Compile: sets loss and optimizer
model.compile(loss='categorical_crossentropy', optimizer='adam')

# Train: forward -> loss -> backward -> update (x epochs)
history = model.fit(X_train, Y_train, epochs=15, batch_size=64)
\end{lstlisting}


% ============================================================
% APPENDIX
% ============================================================
\appendix
\chapter{Derivative Reference Table}

\begin{center}
\renewcommand{\arraystretch}{2.0}
\begin{tabular}{lcc}
\toprule
\textbf{Function} & \textbf{Formula} $f(z)$ & \textbf{Derivative} $f'(z)$ \\
\midrule
ReLU & $\max(0, z)$ & $\begin{cases} 1 & z > 0 \\ 0 & z \leq 0 \end{cases}$ \\[8pt]
Sigmoid & $\dfrac{1}{1+e^{-z}}$ & $\sigma(z)(1-\sigma(z))$ \\[8pt]
Tanh & $\dfrac{e^z - e^{-z}}{e^z + e^{-z}}$ & $1 - \tanh^2(z)$ \\[8pt]
Linear & $z$ & $1$ \\[8pt]
Softmax (elem $i$) & $\dfrac{e^{z_i}}{\sum_j e^{z_j}}$ & $s_i(\delta_{ij} - s_j)$ \\
\bottomrule
\end{tabular}
\end{center}

\chapter{Matrix Calculus Identities}

These identities are used throughout the backpropagation derivations:

\begin{enumerate}
    \item If $\mathbf{z} = W\mathbf{x} + \mathbf{b}$, then $\frac{\partial \mathbf{z}}{\partial W} = \mathbf{x}^\top$ and $\frac{\partial \mathbf{z}}{\partial \mathbf{x}} = W^\top$.
    
    \item If $\mathcal{L} = f(\mathbf{z})$ and $\mathbf{z} = g(\mathbf{x})$, then $\frac{\partial \mathcal{L}}{\partial \mathbf{x}} = \frac{\partial \mathcal{L}}{\partial \mathbf{z}} \cdot \frac{\partial \mathbf{z}}{\partial \mathbf{x}}$.
    
    \item Element-wise operations: If $\mathbf{a} = \sigma(\mathbf{z})$ element-wise, then $\frac{\partial \mathcal{L}}{\partial \mathbf{z}} = \frac{\partial \mathcal{L}}{\partial \mathbf{a}} \odot \sigma'(\mathbf{z})$.
    
    \item For cross-entropy + softmax: $\frac{\partial \mathcal{L}}{\partial \mathbf{z}} = \hat{\mathbf{y}} - \mathbf{y}$ (the elegantly simple combined gradient).
\end{enumerate}

\end{document}
